Our experiments suggest several conclusions that are relevant for further investigation in model-based reinforcement learning. 
First, our results show that model-based reinforcement learning with neural network dynamics models can achieve results that are competitive not only with Bayesian nonparametric models such as GPs, but also on par with model-free algorithms such as PPO and SAC in terms of asymptotic performance, while attaining substantially more efficient convergence. 
Although the individual components of our model-based reinforcement learning algorithms are not individually new -- prior works have suggested both ensembling and outputting Gaussian distribution parameters~\citep{Lakshminarayanan2017}, as well as the use of MPC for model-based RL~\citep{Nagabandi2017} -- the particular combination of these components into a model-based reinforcement learning algorithm is, to our knowledge, novel, and the results provide a new state-of-the-art for model-based reinforcement learning algorithms based on high-capacity parametric models such as neural networks. The systematic investigation in our experiments was a critical ingredient in determining the precise combination of these components that attains the best performance.

Our results indicate that the gap in asymptotic performance between model-based and model-free reinforcement learning can, at least in part, be bridged by incorporating uncertainty estimation into the model learning process. 
Our experiments further indicate that both epistemic and aleatoric uncertainty plays a crucial role in this process. 
Our analysis considers a model-based algorithm based on dynamics estimation and planning. 
A compelling alternative class of methods uses the model to train a parameterized policy~\citep{Ko2007,Deisenroth2014,McAllister2017}. 
While the choice of using the model for planning versus policy learning is largely orthogonal to the other design choices, a promising direction for future work is to investigate how policy learning can be incorporated into our framework to amortize the cost of planning at test-time. 
Our initial experiments with policy learning did not yield an effective algorithm by directly propagating gradients through our uncertainty-aware models. 
We believe this may be due to chaotic policy gradients, whose recent analysis~\citep{Parmas2018} could help yield a policy-based PETS in future work.
Finally, the observation that model-based RL can match the performance of model-free algorithms suggests that substantial further investigation of such of methods is in order, as a potential avenue for effective, sample-efficient, and practical general-purpose reinforcement learning.
